\documentclass[11pt]{amsart}
\usepackage{calc,amssymb,amsmath,ulem,stmaryrd,fullpage}
\usepackage{enumerate}
\usepackage{alltt}
\RequirePackage[dvipsnames,usenames]{color}
\let\mffont=\sf
\normalem
%\input{kmacros3.sty}
\input{xy}
\xyoption{all}
\usepackage{hyperref}
\usepackage{graphicx}
\usepackage[all,cmtip]{xy}
\usepackage{verbatim}
\usepackage[left=0.95in,top=0.95in,right=0.95in,bottom=0.95in]{geometry}
\newcommand{\tauCohomology}{T}
\newcommand{\FCohomology}{S}


\theoremstyle{theorem}
%\newtheorem*{mainthm*}{Main Theorem}

\newcommand{\note}[1]{\marginpar{\sffamily\tiny #1}}

\def\todo#1{\textcolor{Mahogany}%
{\sffamily\footnotesize\newline{\color{Mahogany}\fbox{\parbox{\textwidth-15pt}{#1}}}\newline}}

\renewcommand{\O}{\mathcal O}
\newcommand{\ie}{{\itshape i.e.} }
\newcommand{\roundup}[1]{\lceil #1 \rceil}
\newcommand{\rounddown}[1]{\lfloor #1 \rfloor}
\renewcommand{\to}[1][]{\xrightarrow{\ #1\ }}
\newcommand{\onto}[1][]{\protect{\xrightarrow{\ #1\ }\hspace{-0.8em}\rightarrow}}
\newcommand{\into}[1][]{\lhook \joinrel \xrightarrow{\ #1\ }}
%\newcommand{DBDual}[1]{{\underline{\omega}_{#1}}}
\newcommand{\DBDual}[1]{{{\uuline \omega}{}^{\mydot}_{#1}}}
\newcommand{\Tt}{{\mathfrak{T}}}
%\newcommand{\bn}{{\mathfrak{n}} }
\newcommand{\sol}[1]{ \vskip 6pt{\color{blue} {\bf Solution:} #1} \vskip 6pt}

\newcommand{\bR}{{\mathbb{R}}}
\newcommand{\va}{{\vec{a}}}
\newcommand{\vb}{{\vec{b}}}
\newcommand{\vzero}{{\vec{0}}}
\newcommand{\vi}{{\vec{i}}}
\newcommand{\vj}{{\vec{j}}}
\newcommand{\vk}{{\vec{k}}}
\newcommand{\vh}{{\vec{h}}}
\newcommand{\vr}{{\vec{r}}}
\newcommand{\vu}{{\vec{u}}}
\newcommand{\vv}{{\vec{v}}}
\newcommand{\vw}{{\vec{w}}}
\newcommand{\vx}{{\vec{x}}}
\newcommand{\vy}{{\vec{y}}}
\newcommand{\vz}{{\vec{z}}}
\newcommand{\vT}{{\vec{T}}}
\newcommand{\vN}{{\vec{N}}}
\newcommand{\vF}{{\vec{F}}}
\newcommand{\vG}{{\vec{G}}}
\newcommand{\bF}{\mathbb{F}}
\newcommand{\bQ}{\mathbb{Q}}
\newcommand{\bZ}{\mathbb{Z}}
\newcommand{\bC}{\mathbb{C}}
\newcommand{\Gal}{\textnormal{Gal}}
\newcommand{\Aut}{\textnormal{Aut}}
\newcommand{\Emb}{\textnormal{Emb}}
\newcommand{\Tr}{\textnormal{Tr}}
\newcommand{\N}{\textnormal{N}}
\newcommand{\Hom}{\textnormal{Hom}}
\newcommand{\Spec}{\textnormal{Spec}}


\begin{document}
\title{HW \#5 -- Math 6320\\Spring 2015}
\author{Due: Tuesday March 10th}
\maketitle

\begin{enumerate}
\item Suppose that $W_1 \subseteq W_2$ are multiplicative sets in a ring $R$ and $M$ is an $R$-module.  Show that $W_2^{-1}(W_1^{-1} M) \cong W_2^{-1} M$.
\item Suppose that $N \subseteq M$ are $R$-modules and $W \subseteq R$ is a multiplicative set.  Show directly that $W^{-1}(M/N) \cong W^{-1} M / W^{-1}N$. Here we are using the fact that $W^{-1} N$ can be identified with a submodule of $W^{-1} M$.
\item  We consider localizing at prime ideals.
   \begin{itemize}
       \item[(a)]  Suppose that $A$ is a ring and that for each prime ideal $P \in \Spec A$, the local ring $A_P := (A \setminus P)^{-1} A$ is an integral domain.  Show that $A$ need not be an integral domain.
       \item[(b)]  However, if $M$ and $N$ are $A$-modules with a map $f : M \to N$.  Consider the induced map $f_P : (A \setminus P)^{-1} M =: M_P \to N_P := (A \setminus P)^{-1} N$ for each $P \in \Spec A$.  Show that $f$ is surjective (respectively injective) if and only if $f_P$ is surjective (respectively injective) for every $P \in \Spec A$.
       \end{itemize}
   \item Suppose that $R$ is an integral domain and $0 \neq r \in R$.  Set $W = \{1, r, r^2, r^3, \ldots\}$.  Is $R[r^{-1}] \cong W^{-1}R \cong R[x]/\langle xr - 1 \rangle$ where $x$ is an indeterminant?
   \item Suppose that $k$ is a field.  Consider the following set
\[
S := \{ f \in k[x] \;|\; f(0) = f(1) \}.
\]
Show that $S$ is a ring.
Describe $S$ as a quotient of a polynomial ring and describe its prime Spectrum.  Additionally, describe the induced map on $\Spec$s from the inclusion of rings $S \hookrightarrow k[x]$ at least in the case that $k$ is algebraically closed.
\item Suppose that $k$ is a field and that $f : R \to S$ is a map between finitely generated $k$-algebras (this means that $R$ is of the form $k[x_1, \dots, x_n]/I$, and likewise with $S$, and also that $f$ sends $k$ to $k$).  Show that the function $f^\# : \Spec S \to \Spec R$ sends maximal ideals to maximal ideals.  
    \vskip 3pt
\emph{Hint:} You may use the fact that if $K$ is a field and $L \supseteq K$ a field extension such that $L$ is a finitely generated $K$-algebra, then $L$ is a finite field extension.  You may also use the fact that an integral domain which is a module finite extension of a field is itself a field.
\item Suppose that $A$ and $B$ are rings.  Figure out what $\Spec (A \oplus B)$ is in relation to $\Spec A$ and $\Spec B$.
\item (\emph{Warning: Hard}) Suppose that $A$ and $B$ are rings, $I$ is an ideal of $A$ with canonical projection $f : A \to A/I$.  Further suppose that we are given a ring homomorphism $g : B \to A/I$.  Consider the following set.
\[
C = \{ (a,b) \in A \oplus B | f(a) = g(b) \}.
\]
\begin{itemize}
\item[(a)]  Show that $C$ is a subring of $A \oplus B$.  Note that $C$ has canoical maps to $A$ and to $B$ (call them $p_1$ and $p_2$ respectively).
\item[(b)]  Show that the elements of $\Spec C$ are in bijection with the set
\[
\Big( (\Spec A) \setminus V(I) \Big) \coprod (\Spec B).
\]
\emph{Hint: } Consider a prime in $\Spec C$.  There are two possibilities, it either contains $\ker p_2$ or it does not.  In the latter case, invert an appropriate element and analyze what happens.
\item[(c)]  Describe geometrically $\Spec C$ in the following examples where $k$ is an algebraically closed field.  Additionally describe the map $\Spec A \to \Spec C$ in each case:
\begin{itemize}
\item[(i)]  $A := k[x]$, $I = \langle x^2 - 1 \rangle$ and $B = k$.
\item[(ii)]  $A := k[x,y]$, $I = \langle x \rangle$ and $B = k$.
\end{itemize}
\item[(d)]  Suppose now that $g$ is surjective.  Explain why $\Spec C$ is the gluing of $\Spec A$ to $\Spec B$ along $V(I)$.
\end{itemize}
The following philosophical statement on the elements of $C$ might help with this problem.  $C$ is made up of the functions in $A \oplus B$ that agree on the set $V(I)$.
\end{enumerate}

\end{document}

